% Ampliación demostrativa DF003--DF007. No modifica la fuente integral.
% Propietario común de los archivos Lean citados en estos comentarios:
% output/PAQUETE_ARTICULO_I_INTEGRACION_ESTADO_CAMPOS_20260919/.
% Residencia causal de la primera sección: después de los lectores regionales
% coinductivos y antes de emplear sus registros como entrada de la selección K.
% Residencia causal de la segunda sección: después de la construcción marcada
% Paley--Witt--Golay--Leech y antes de la reconstrucción de la estructura VOA.
% Los dos cortes expositivos conservan el mismo origen APP--TRIT--TPK.

\section[Generación regional y refinamiento]{Compatibilidad integral de la generación regional, el refinamiento cilíndrico y la incidencia}
\label{act21:sec:df-lectores-cilindros}

La prolongación regional produce sucesiones ordenadas de bloques, junto con
sus profundidades de lectura. La proyección arquimediana y la incidencia
ternaria se calculan sobre esas mismas sucesiones. El vínculo entre ambas
realizaciones puede formularse enteramente mediante operaciones enteras:
selección por intervalos racionales, división euclídea, elevación de seis
trits, lectura de acarreo y transformación de Witt. Esta formulación permite
seguir cada coeficiente desde su bloque generador hasta su empleo posterior.

La construcción utiliza los tres lectores regionales ya obtenidos del estado
APP--TRIT--TPK. Los denotaremos por los índices
\(c\in\{\mathrm P,\mathrm E,\mathrm F\}\), correspondientes, respectivamente,
a las regiones de clausura, propagación y autoescala. Su realización real se
escribirá \(v_c\). Las identificaciones posteriores
\(v_{\mathrm P}=\pi_{\mathrm{HMT}}\),
\(v_{\mathrm E}=e_{\mathrm{HMT}}\) y
\(v_{\mathrm F}=\varphi_{\mathrm{HMT}}\) conservan ese orden causal.
La operación que emite un bloque inspecciona los intervalos racionales del
lector regional; el valor real interviene en la demostración de corrección
de dicha operación.

\subsection{Emisión posicional mediante separación racional de celdas}
\label{act21:subsec:df-emision-racional}

% Fuente: lean/biblioteca/RegionalPublicationComposition.lean,
% produced_brackets, eventual_publication, joint_eventual_publication,
% search_terminates, stopped_cell_correct, publications_compatible.

Para cada región se dispone de dos sucesiones racionales
\(\ell_c(j),u_c(j)\), construidas por su lector, que cumplen
\[
 \ell_c(j)\leq v_c\leq u_c(j).
\]
La propiedad de separación de celdas demostrada para esos lectores afirma
que, para cada entero \(s>0\), existe \(J\) tal que, para todo \(j\geq J\),
\begin{equation}
 \lfloor s\ell_c(j)\rfloor
 =\lceil su_c(j)\rceil-1
 =\lfloor sv_c\rfloor.
 \label{act21:eq:df-separacion-celdas}
\end{equation}
La evitación de las fronteras racionales por los valores regionales es la
hipótesis de los teoremas previos que garantiza esta separación. Aquí se
emplea su conclusión operativa, con los lectores y los intervalos que la
producen.

\begin{definition}[Índice de separación y prefijo emitido]
Para \(s>0\), sea
\[
 j_c(s)=\min\{j\in\mathbb N:
 \lfloor s\ell_c(j)\rfloor=\lceil su_c(j)\rceil-1\}.
\]
Para una base entera \(b\geq2\) y una profundidad \(n\geq0\), se define
\begin{equation}
 P_c(b,n)=\lfloor b^n\ell_c(j_c(b^n))\rfloor.
 \label{act21:eq:df-publicacion-prefijo}
\end{equation}
El prefijo contiene la parte entera y las primeras \(n\) posiciones en base
\(b\); por consiguiente, \(P_c(b,0)\) es la parte entera regional.
\end{definition}

\begin{theorem}[Terminación y corrección de la emisión]
\label{act21:thm:df-emision-correcta}
El índice \(j_c(s)\) existe para todo \(s>0\). Los prefijos emitidos cumplen
\begin{equation}
 P_c(b,n)=\lfloor b^n v_c\rfloor,
 \qquad
 P_c(b,n)\leq b^n v_c<P_c(b,n)+1.
 \label{act21:eq:df-prefijo-correcto}
\end{equation}
Para cada escala \(s\) puede elegirse además un único índice de suficiencia
\(J(s)\) válido simultáneamente para las tres regiones.
\end{theorem}

\begin{proof}
La igualdad eventual de \eqref{act21:eq:df-separacion-celdas} hace no vacío el
conjunto que define \(j_c(s)\). Su mínimo existe por buen orden de
\(\mathbb N\). En cualquier índice donde se produzca la igualdad, la cota
inferior racional proporciona
\(\lfloor s\ell_c(j)\rfloor\leq\lfloor sv_c\rfloor\).
La cota superior, junto con la evitación de una frontera de celda, da
\(\lfloor sv_c\rfloor\leq\lceil su_c(j)\rceil-1\).
Al coincidir ambos extremos enteros, el entero intermedio coincide con
ellos. Esto demuestra \eqref{act21:eq:df-prefijo-correcto}. Para la afirmación
conjunta basta tomar el máximo de los tres índices de suficiencia obtenidos
para la misma escala. El índice conjunto controla simultáneamente las tres
emisiones, mientras que cada región conserva sus propios intervalos.
\end{proof}

\begin{proposition}[Compatibilidad exacta de los prefijos]
\label{act21:prop:df-prefijos-compatibles}
Para \(n,m\geq0\),
\begin{equation}
 \left\lfloor\frac{P_c(b,n+m)}{b^m}\right\rfloor=P_c(b,n).
 \label{act21:eq:df-truncamiento}
\end{equation}
En particular, el bloque siguiente
\(d_c(b,n)=P_c(b,n+1)\bmod b\) satisface
\begin{equation}
 P_c(b,n+1)=bP_c(b,n)+d_c(b,n),\qquad 0\leq d_c(b,n)<b.
 \label{act21:eq:df-paso-posicional}
\end{equation}
\end{proposition}

\begin{proof}
Escribamos \(b^n v_c=N+\theta\), con
\(N=P_c(b,n)\) y \(0\leq\theta<1\). Entonces
\(P_c(b,n+m)=b^mN+\lfloor b^m\theta\rfloor\), y el segundo sumando
pertenece a \(\{0,\ldots,b^m-1\}\). La división euclídea por \(b^m\)
recupera \(N\). La especialización \(m=1\) produce
\eqref{act21:eq:df-paso-posicional} y su cota residual.
\end{proof}

\subsection{Bloques regionales de seis trits a profundidad arbitraria}
\label{act21:subsec:df-bloques-arbitrarios}

% Fuente: deltas/integracion/JointRegionalFrontier.lean,
% publish_base729_eq_base3, generated_block_from_longer_prefix,
% generated_block_eq_finite, generated_signature_eq_finite.

La igualdad entera \(729=3^6\) determina la agrupación de seis posiciones
ternarias en un bloque. Para \(t\geq0\), definimos
\begin{equation}
 u_c(t)=P_c(729,t+1)\bmod729,
 \qquad
 B_{c,j}(t)=
 \left\lfloor\frac{u_c(t)}{3^{5-j}}\right\rfloor\bmod3,
 \quad 0\leq j<6.
 \label{act21:eq:df-bloque-y-banda}
\end{equation}
De este modo, \(B_c(t)\) es una palabra ordenada de seis trits y
\begin{equation}
 u_c(t)=\sum_{j=0}^{5}B_{c,j}(t)3^{5-j}.
 \label{act21:eq:df-elevacion-banda}
\end{equation}
La palabra y su elevación entera son dos representaciones mutuamente
recuperables del mismo bloque emitido.

\begin{theorem}[Estabilidad de extracción a cualquier profundidad]
\label{act21:thm:df-bloque-estable}
Se cumplen, para todo \(n\),
\[
 P_c(729,n)=P_c(3,6n).
\]
Si \(t<n\), entonces
\begin{equation}
 u_c(t)=
 \left\lfloor\frac{P_c(729,n)}{729^{n-t-1}}\right\rfloor\bmod729.
 \label{act21:eq:df-bloque-prefijo-largo}
\end{equation}
En consecuencia, cualquier ampliación del prefijo preserva todos los
bloques previamente emitidos y todas las firmas que se calculan sobre ellos.
\end{theorem}

\begin{proof}
La primera igualdad resulta de que ambos prefijos se separan a la misma
escala, \(729^n=3^{6n}\), y del
teorema~\ref{act21:thm:df-emision-correcta}. Para la segunda aplicamos
\eqref{act21:eq:df-truncamiento} a las profundidades \(n\) y \(t+1\), y después
tomamos el residuo módulo \(729\). La extracción ternaria de
\eqref{act21:eq:df-bloque-y-banda} es determinista. Cualquier función explícita
de esas dieciocho coordenadas conserva, por tanto, el mismo resultado al
ampliar el prefijo.
\end{proof}

La construcción está cuantificada sobre \(t\in\mathbb N\). Una tabla que
materialice cien bloques constituye una evaluación finita de esta
construcción; la fórmula \eqref{act21:eq:df-bloque-prefijo-largo} determina
exactamente su compatibilidad con cualquier extensión posterior.

\begin{example}[Primer bloque conjunto]
La evaluación de los lectores regionales seleccionados da las palabras
\[
 B_{\mathrm P}(0)=010211,\qquad
 B_{\mathrm E}(0)=201101,\qquad
 B_{\mathrm F}(0)=121200.
\]
Su elevación por \eqref{act21:eq:df-elevacion-banda} produce
\[
 u_{\mathrm P}(0)=103,\qquad
 u_{\mathrm E}(0)=523,\qquad
 u_{\mathrm F}(0)=450.
\]
Por ejemplo,
\(201101_3=2\cdot243+27+9+1=523\).
Las seis posiciones, incluidos los ceros iniciales, se conservan en la
representación de banda. El cero inicial de \(010211\) tiene una posición
definida aunque la escritura decimal de \(103\) carezca de él.
\end{example}

\subsection{Registros de transición y reconstrucción de las elevaciones lineales}
\label{act21:subsec:df-registros-transicion}

% Fuentes: deltas/integracion/RegionalTransitionRecords.lean y
% GeneratedTransitionRecords.lean; lean/biblioteca/TPKLifts.lean.
% Los enlaces reconstruidos son R(0)->R(1) y R(2)->R(3).

Sobre \(\mathbb F_3\), reunimos dos emisiones consecutivas de cada región
en la matriz
\begin{equation}
 R(t)=
 \begin{pmatrix}
 B_{\mathrm P}(t)\\ B_{\mathrm P}(t+1)\\
 B_{\mathrm E}(t)\\ B_{\mathrm E}(t+1)\\
 B_{\mathrm F}(t)\\ B_{\mathrm F}(t+1)
 \end{pmatrix}.
 \label{act21:eq:df-registro-matricial}
\end{equation}
Esta definición fija tanto el orden regional como el orden temporal. Los
cuatro primeros registros son
\[
 X_0=R(0),\quad Y_0=R(1),\quad X_1=R(2),\quad Y_1=R(3).
\]
Sus valores, obtenidos por extracción de los prefijos regionales, son
\begin{align*}
 X_0&=\begin{pmatrix}
0&1&0&2&1&1\\0&1&2&2&2&2\\2&0&1&1&0&1\\
1&2&1&2&2&1\\1&2&1&2&0&0\\1&1&2&2&0&2
\end{pmatrix},&
 Y_0&=\begin{pmatrix}
0&1&2&2&2&2\\0&1&0&2&1&1\\1&2&1&2&2&1\\
1&0&2&0&1&1\\1&1&2&2&0&2\\1&2&1&0&2&0
\end{pmatrix},\\[1ex]
 X_1&=\begin{pmatrix}
0&1&0&2&1&1\\0&0&2&1&1&1\\1&0&2&0&1&1\\
0&1&2&2&2&2\\1&2&1&0&2&0\\0&1&0&2&1&0
\end{pmatrix},&
 Y_1&=\begin{pmatrix}
0&0&2&1&1&1\\1&1&0&2&2&1\\0&1&2&2&2&2\\
1&0&2&0&1&1\\0&1&0&2&1&0\\0&1&0&2&0&0
\end{pmatrix}.
\end{align*}

\begin{theorem}[Reconstrucción única de los dos enlaces iniciales]
\label{act21:thm:df-elevaciones-reconstruidas}
Las ecuaciones \(X_0L_0=Y_0\) y \(X_1L_1=Y_1\) tienen, cada una, una
única solución en \(M_6(\mathbb F_3)\). Estas soluciones son
\begin{align*}
 L_0&=\begin{pmatrix}
2&2&2&1&2&1\\2&1&2&2&1&1\\1&0&0&1&0&2\\
0&0&1&1&1&0\\2&2&2&0&2&1\\2&1&2&1&0&0
\end{pmatrix},&
 L_1&=\begin{pmatrix}
2&1&2&0&2&2\\1&2&1&1&2&0\\1&2&1&1&1&2\\
0&1&0&0&2&1\\2&0&2&1&0&1\\0&2&2&2&1&1
\end{pmatrix}.
\end{align*}
Ambas matrices son invertibles.
\end{theorem}

\begin{proof}
Las matrices
\begin{align*}
 A_0&=\begin{pmatrix}
1&1&2&0&1&2\\0&0&1&0&0&1\\2&1&2&1&0&1\\
0&2&0&1&0&2\\2&1&1&0&2&2\\2&1&1&1&1&2
\end{pmatrix},&
 A_1&=\begin{pmatrix}
1&0&1&2&0&0\\0&0&1&1&2&2\\0&1&2&0&1&1\\
1&1&0&2&0&0\\1&1&2&1&1&2\\1&0&0&0&0&2
\end{pmatrix}
\end{align*}
satisfacen \(A_iX_i=X_iA_i=I\), por multiplicación en \(\mathbb F_3\).
Por tanto, las soluciones están forzadas por \(L_i=A_iY_i\); la
multiplicación produce las dos matrices mostradas. Sus inversas son
\begin{align*}
 L_0^{-1}&=\begin{pmatrix}
1&2&0&2&0&2\\2&0&2&2&0&0\\2&1&2&0&2&0\\
1&0&0&0&2&0\\0&2&1&1&2&0\\2&2&2&2&2&2
\end{pmatrix},&
 L_1^{-1}&=\begin{pmatrix}
2&0&2&1&2&1\\2&1&0&0&2&0\\2&0&2&2&0&2\\
2&0&0&1&1&0\\1&1&1&1&1&0\\2&0&1&2&2&0
\end{pmatrix}.
\end{align*}
La verificación bilateral \(L_iL_i^{-1}=L_i^{-1}L_i=I\) demuestra la
recuperación exacta de cada palabra transformada. Todos los productos
utilizan residuos \(0,1,2\), con su interpretación en \(\mathbb F_3\).
\end{proof}

La lectura causal de este resultado es precisa: las emisiones regionales
producen \(R(t)\); los registros \(R(0),\ldots,R(3)\) determinan las dos
elevaciones anteriores; sus inversas recuperan las palabras transformadas.
La continuidad para todo \(t\) pertenece a los lectores coinductivos y al
teorema~\ref{act21:thm:df-bloque-estable}. La reconstrucción matricial aquí
demostrada tiene por dominio los dos enlaces especificados en su enunciado.

\subsection{Refinamiento simultáneo en las bases ternaria y decimal}
\label{act21:subsec:df-cilindros-mixtos}

% Fuente: deltas/integracion/RegionalCylinderTransition.lean.

Una emisión ternaria de seis trits y una emisión decimal de tres cifras
constituyen dos lecturas de distinta base. Para conservar su relación se
registra el estado
\[
 s=(n,k,N,D),
\]
donde \(N\) es un prefijo a profundidad \(n\) en base \(729\), y \(D\) es
un prefijo a profundidad \(k\) en base \(1000\). Sus cilindros
arquimedianos son
\begin{equation}
 I_{729}(s)=\left[\frac{N}{729^n},\frac{N+1}{729^n}\right),\qquad
 I_{1000}(s)=\left[\frac{D}{1000^k},\frac{D+1}{1000^k}\right).
 \label{act21:eq:df-cilindros}
\end{equation}
El extremo superior semiabierto asigna cada valor a una única celda de una
base y una profundidad dadas.

\begin{definition}[Defectos enteros de compatibilidad]
Definimos
\begin{align}
 A(s)&=N1000^k-D729^n,\label{act21:eq:df-defecto-inferior}\\
 B(s)&=(N+1)1000^k-D729^n=A(s)+1000^k.
 \label{act21:eq:df-defecto-superior}
\end{align}
El estado es compatible cuando
\begin{equation}
 A(s)<729^n,\qquad B(s)>0.
 \label{act21:eq:df-compatibilidad}
\end{equation}
\end{definition}

\begin{lemma}[Criterio entero de intersección]
Los cilindros de \eqref{act21:eq:df-cilindros} tienen intersección no vacía si y
sólo si se cumple \eqref{act21:eq:df-compatibilidad}.
\end{lemma}

\begin{proof}
Dos intervalos semiabiertos \([a,b)\) y \([c,d)\), con \(a<b\) y
\(c<d\), se cortan exactamente cuando \(a<d\) y \(c<b\). Aplicando esta
condición a \eqref{act21:eq:df-cilindros} y multiplicando por los denominadores
positivos, obtenemos
\(N1000^k<(D+1)729^n\) y
\(D729^n<(N+1)1000^k\). Son las dos desigualdades de
\eqref{act21:eq:df-compatibilidad}.
\end{proof}

\begin{definition}[Actualización de profundidades y prefijos]
Para \(0\leq u<729\) y \(0\leq d<1000\), definimos
\begin{align*}
 T_u(n,k,N,D)&=(n+1,k,729N+u,D),\\
 T_{u,d}(n,k,N,D)&=(n+1,k+1,729N+u,1000D+d).
\end{align*}
La primera operación refina únicamente el cilindro ternario. La segunda
refina ambas lecturas y conserva las dos profundidades en el estado.
\end{definition}

\begin{proposition}[Recurrencias exactas de frontera]
\label{act21:prop:df-rec-frontera}
Las actualizaciones anteriores satisfacen
\begin{align}
 A(T_u s)&=729A(s)+u1000^k,\label{act21:eq:df-rec-ternaria}\\
 A(T_{u,d}s)&=729000A(s)+u1000^{k+1}-d729^{n+1}.
 \label{act21:eq:df-rec-conjunta}
\end{align}
\end{proposition}

\begin{proof}
Para la primera identidad,
\[
 (729N+u)1000^k-D729^{n+1}
 =729(N1000^k-D729^n)+u1000^k.
\]
Para la segunda,
\begin{align*}
 &(729N+u)1000^{k+1}-(1000D+d)729^{n+1}\\
 &\qquad=729000(N1000^k-D729^n)
      +u1000^{k+1}-d729^{n+1}.
\end{align*}
Ambas son identidades enteras anteriores a cualquier aproximación real.
\end{proof}

\begin{theorem}[Refinamiento de los cilindros regionales generados]
\label{act21:thm:df-refinamiento-generado}
Sea
\[
 s_c(n,k)=(n,k,P_c(729,n),P_c(1000,k)).
\]
Entonces \(s_c(n,k)\) es compatible para cualesquiera \(n,k\). Además,
\begin{align*}
 T_{u_c(n)}s_c(n,k)&=s_c(n+1,k),\\
 T_{u_c(n),d_c(1000,k)}s_c(n,k)&=s_c(n+1,k+1).
\end{align*}
La división euclídea de los prefijos refinados recupera exactamente los
prefijos anteriores. Para todo \(a,b\geq0\),
\[
 \frac{P_c(729,n+a)-P_c(729,n+a)\bmod729^a}{729^a}
 =P_c(729,n),
\]
y se cumple la identidad análoga en base \(1000\).
\end{theorem}

\begin{proof}
El teorema~\ref{act21:thm:df-emision-correcta} sitúa \(v_c\) simultáneamente en
los dos intervalos \eqref{act21:eq:df-cilindros}; el criterio de intersección da
la compatibilidad. Las identidades de actualización son
\eqref{act21:eq:df-paso-posicional} en cada base. La recuperación se obtiene de
\eqref{act21:eq:df-truncamiento}. Por tanto, la recurrencia del defecto y la
recuperación del prefijo son consecuencias de una misma actualización.
\end{proof}

\begin{example}[Primer par de cilindros regionales]
Para \(n=k=1\), la evaluación de las tres regiones da
\[
\begin{array}{c|rr|rr}
 c&N&D&A&B\\\hline
 \mathrm P&2290&3141&211&1211\\
 \mathrm E&1981&2718&-422&578\\
 \mathrm F&1179&1618&-522&478
\end{array}
\]
En las tres filas \(A<729\) y \(B>0\). La primera columna de prefijos
se descompone como \(2290=3\cdot729+103\),
\(1981=2\cdot729+523\) y \(1179=729+450\): reaparecen exactamente los
bloques de seis trits ya emitidos. Los signos diferentes de \(A\) expresan
la posición relativa de los dos extremos inferiores, sin alterar la
pertenencia común del valor regional.
\end{example}

\subsection{Selección del bloque siguiente mediante el lector regional}

La compatibilidad de dos intervalos expresa una relación entre lecturas.
La determinación del siguiente bloque utiliza además el intervalo racional
producido por la región. Esta operación puede formularse sin eliminar la
historia de profundidad.

\begin{definition}[Admisibilidad de la prolongación regional]
Fijados \(c,n\), sea \(s=729^{n+1}\) y \(j=j_c(s)\). Un residuo
\(u\in\{0,\ldots,728\}\) es admisible cuando
\begin{equation}
 \lfloor s\ell_c(j)\rfloor
 \leq729P_c(729,n)+u
 \leq\lceil su_c(j)\rceil-1.
 \label{act21:eq:df-admisibilidad}
\end{equation}
\end{definition}

\begin{theorem}[Unicidad de la prolongación compatible con el lector]
\label{act21:thm:df-prolongacion-unica}
La condición \eqref{act21:eq:df-admisibilidad} equivale a \(u=u_c(n)\).
Por tanto, existe exactamente una prolongación admisible de cada región
a cada profundidad.
\end{theorem}

\begin{proof}
En el índice de separación, los dos extremos enteros de
\eqref{act21:eq:df-admisibilidad} coinciden con \(P_c(729,n+1)\).
La condición equivale así a
\(729P_c(729,n)+u=P_c(729,n+1)\).
Por \eqref{act21:eq:df-paso-posicional}, el lado derecho es
\(729P_c(729,n)+u_c(n)\). La cancelación del término común da la
igualdad de residuos. Recíprocamente, esa igualdad satisface la condición.
\end{proof}

La dependencia efectiva de esta selección queda completamente determinada:
región seleccionada, lector racional, profundidad, índice de separación,
prefijo anterior y residuo siguiente. El defecto \(A\) resume una relación
de frontera; el estado \((n,k,N,D)\) y los intervalos regionales conservan
los datos empleados por la actualización.

\subsection{Firma ternaria, acarreo, incidencia de Witt y reloj de lectura}
\label{act21:subsec:df-firma-conjunta}

% Fuentes: lean/regional/N69RegionalSignature.lean,
% deltas/integracion/JointRegionalFrontier.lean y JointCylinderSignature.lean.

Fijada una profundidad \(t\), escribamos \(B_{c,j}=B_{c,j}(t)\). Para
cada columna \(j\) se calculan los siguientes enteros y residuos:
\begin{align}
 S_j&=B_{\mathrm P,j}+B_{\mathrm E,j}+B_{\mathrm F,j},
 &q_j&=S_j\bmod3,
 &c_j&=\left\lfloor S_j/3\right\rfloor,
 \label{act21:eq:df-firma-qc}\\
 a_j&=(B_{\mathrm P,j}+B_{\mathrm E,j}-B_{\mathrm F,j})\bmod3,
 &w_j&=\sum_c\mathbf1_{B_{c,j}\ne0},
 &\bar w_j&=w_j\bmod3.
 \label{act21:eq:df-firma-axial}
\end{align}
El residuo de la expresión axial se toma en los enteros. Una implementación
con naturales lo calcula como
\((B_{\mathrm P,j}+B_{\mathrm E,j}+3-B_{\mathrm F,j})\bmod3\), que conserva
exactamente ese residuo.

La matriz de incidencia ternaria utilizada por el lector es
\begin{equation}
 W=\begin{pmatrix}
0&1&1&1&1&1\\
1&0&1&1&2&2\\
1&1&0&2&1&2\\
1&1&2&0&2&1\\
1&2&1&2&0&1\\
1&2&2&1&1&0
\end{pmatrix}.
\label{act21:eq:df-matriz-witt}
\end{equation}
Para cada región se conservan dos cargas residuales:
\begin{equation}
 v_c^{\mathrm{res}}=\sum_j B_{c,j}\bmod3,
 \qquad
 v_c^{\mathrm W}=\sum_j(B_cW)_j\bmod3.
 \label{act21:eq:df-cargas-residuales}
\end{equation}
La firma completa del lector es
\(\mathcal S(B)=(q,a,c,w,\bar w,v^{\mathrm{res}},v^{\mathrm W})\).
Su estructura conserva por separado el residuo de columna, el acarreo, el
contraste axial, el peso de soporte y las dos lecturas de carga.

\begin{theorem}[Recuperación exacta y cotas de la firma]
\label{act21:thm:df-recuperacion-firma}
Para cada columna y cada región se cumplen
\begin{align}
 S_j&=q_j+3c_j,&0\leq q_j,a_j&<3,&0\leq c_j&\leq2,
 \label{act21:eq:df-recuperacion-total}\\
 B_{\mathrm F,j}&=2(q_j+3-a_j)\bmod3,&0\leq w_j&\leq3,
 &0\leq\bar w_j&<3.
 \label{act21:eq:df-recuperacion-eje}
\end{align}
Además,
\begin{equation}
 v_c^{\mathrm W}
 =\left(\sum_j\bigl((B_cW)_j\bmod3\bigr)\right)\bmod3.
 \label{act21:eq:df-witt-reduccion}
\end{equation}
Las recuperaciones enunciadas determinan el total de cada columna y su
coordenada de autoescala; el registro regional conserva adicionalmente
las tres bandas ordenadas.
\end{theorem}

\begin{proof}
La identidad \eqref{act21:eq:df-recuperacion-total} es la división euclídea de
\(S_j\) por \(3\). Cada sumando de \(S_j\) está entre \(0\) y \(2\),
de modo que \(0\leq S_j\leq6\) y \(c_j\leq2\). En \(\mathbb F_3\),
la resta de las dos congruencias que definen \(q_j\) y \(a_j\) da
\(q_j-a_j=2B_{\mathrm F,j}\). Como el inverso de \(2\) es \(2\), se
obtiene \eqref{act21:eq:df-recuperacion-eje}; el sumando \(3\) permite
escribirla con naturales. Las cotas de \(w_j\) resultan de sumar tres
indicadores. Finalmente, reducir una suma entera módulo \(3\) equivale
a reducir cada sumando y volver a reducir la suma, lo que prueba
\eqref{act21:eq:df-witt-reduccion}.
\end{proof}

\begin{example}[Firma del primer bloque conjunto]
Las tres bandas iniciales producen
\begin{align*}
 q&=(0,0,2,2,1,2),&a&=(1,2,0,1,1,2),\\
 c&=(1,1,0,1,0,0),&w&=(2,2,2,3,1,2),\\
 v^{\mathrm{res}}&=(2,2,0),&v^{\mathrm W}&=(2,1,1).
\end{align*}
La cuarta columna tiene total \(2+1+2=5\); su firma registra
\(q_3=2\) y \(c_3=1\), por lo que \(q_3+3c_3=5\).
La recuperación axial da
\(2(2+3-1)\bmod3=2\), precisamente la coordenada de autoescala.
Las imágenes de Witt, reducidas componente a componente, son
\[
 (2,0,2,1,1,2),\qquad
 (0,0,0,2,0,2),\qquad
 (2,1,1,2,1,0).
\]
Su suma residual por fila produce las tres cargas \((2,1,1)\).
\end{example}

\begin{definition}[Reloj entero de lectura decimal]
El índice de transición \(t\) y la profundidad decimal de representación se
relacionan mediante
\begin{equation}
 \kappa(t)=\max\{k\in\mathbb N:1000^k\leq3^{6(t+1)}\}.
 \label{act21:eq:df-reloj}
\end{equation}
\end{definition}

\begin{proposition}[Caracterización y monotonía del reloj]
\label{act21:prop:df-reloj}
El entero \(\kappa(t)\) queda determinado unívocamente por
\begin{equation}
 1000^{\kappa(t)}\leq729^{t+1}<1000^{\kappa(t)+1}.
 \label{act21:eq:df-reloj-cotas}
\end{equation}
La función \(\kappa\) es monótona y satisface
\(\kappa(20)=\kappa(21)=20\).
\end{proposition}

\begin{proof}
Las potencias de \(1000\) son estrictamente crecientes y no acotadas;
existe, por tanto, un único intervalo consecutivo de potencias que
contiene \(729^{t+1}\). La monotonía de esta última expresión demuestra
la de \(\kappa\). Las dos evaluaciones finales son las comparaciones
enteras
\[
 1000^{20}\leq729^{21}<1000^{21},\qquad
 1000^{20}\leq729^{22}<1000^{21}.
\]
Estas comparaciones se efectúan con potencias enteras exactas.
\end{proof}

La igualdad de profundidad decimal en dos transiciones distintas explica
la necesidad de conservar ambos índices: el reloj de lectura puede
permanecer constante mientras avanza la banda ternaria y se actualiza su
firma. La información cronológica se encuentra en el registro ordenado
de las transiciones.

\subsection{Teorema de composición del bloque, el cilindro y la firma}

\begin{theorem}[Compatibilidad conjunta de las lecturas regionales]
\label{act21:thm:df-composicion-conjunta}
Para cada profundidad \(n\), existe una única terna de bloques
\((u_{\mathrm P},u_{\mathrm E},u_{\mathrm F})\) que satisface la
admisibilidad regional \eqref{act21:eq:df-admisibilidad} en las tres regiones.
Esta terna es \((u_{\mathrm P}(n),u_{\mathrm E}(n),u_{\mathrm F}(n))\).
Simultáneamente:
\begin{enumerate}
\item cada bloque actualiza su cilindro por las recurrencias
\eqref{act21:eq:df-rec-ternaria} y \eqref{act21:eq:df-rec-conjunta};
\item sus seis trits producen la firma
\(\mathcal S(B_{\mathrm P}(n),B_{\mathrm E}(n),B_{\mathrm F}(n))\);
\item los totales de columna se recuperan como \(q+3c\), y las cargas
de Witt se obtienen de las mismas bandas mediante
\eqref{act21:eq:df-witt-reduccion};
\item todo prefijo de mayor profundidad devuelve exactamente estos
bloques, cilindros truncados y firmas.
\end{enumerate}
\end{theorem}

\begin{proof}
Aplicamos el teorema~\ref{act21:thm:df-prolongacion-unica} a cada uno de los tres
índices regionales. La unicidad coordenada a coordenada da la unicidad
conjunta. Sustituimos esos mismos residuos en las operaciones
\(T_u,T_{u,d}\), lo que permite aplicar el
teorema~\ref{act21:thm:df-refinamiento-generado}. La aplicación determinista de
\eqref{act21:eq:df-bloque-y-banda} produce las bandas utilizadas por
\eqref{act21:eq:df-firma-qc}--\eqref{act21:eq:df-cargas-residuales}. Las identidades
de recuperación son las del teorema~\ref{act21:thm:df-recuperacion-firma}.
Finalmente, el teorema~\ref{act21:thm:df-bloque-estable} y el truncamiento de
cada base demuestran la estabilidad a mayor profundidad. En cada paso se
ha conservado la identidad del bloque emitido; ninguna de las dos
lecturas requiere seleccionar otro bloque.
\end{proof}

La conclusión es una compatibilidad operacional entre la contracción
cilíndrica y la conservación de incidencia. Ambas utilizan el mismo
registro generado, mantienen las profundidades y admiten recuperación
entera. Este resultado es la forma explícita del enlace
\[
 \text{bloque regional emitido}
 \longrightarrow
 \begin{cases}
 \text{cilindro y frontera entera},\\
 \text{banda, firma, acarreo y carga de Witt}.
 \end{cases}
\]
El desarrollo posterior de la incidencia y de los campos emplea la
segunda lectura, conservando la procedencia común con la primera.
